\section{Neural Network Intelligence (NNI)}
\textbf{Neural Network Intelligence (NNI)}, developed by Microsoft, is an open-source AutoML toolkit designed to facilitate the complex process of model optimization. Within the scope of this thesis, NNI serves as the primary engine for \textit{Hardware-Constraint Optimization}, bridging the gap between high-level architectural search and the rigid physical limits of STM32 microcontrollers. NNI’s modular design allows it to orchestrate multiple experiments across local and distributed environments, leveraging a "Tuner-Assessor" architecture to find the optimal balance between model accuracy and resource consumption \cite{nni_microsoft}.

\subsection{Hyperparameter Optimization and Architectural Search}
The core utility of NNI lies in its ability to automate the selection of hyperparameter sets that would otherwise require hundreds of hours of manual experimentation. In our framework, the agent defines a \textbf{Search Space} encompassing parameters such as layer widths, kernel sizes, and learning rates. NNI then employs advanced optimization algorithms to explore this space:
\begin{itemize}
    \item \textbf{Random and Grid Search:} Basic strategies for broad exploration of the parameter landscape.
    \item \textbf{Bayesian Optimization (TPE):} An advanced strategy that builds a probabilistic model of the objective function (e.g., maximizing accuracy while staying under 256KB RAM) to suggest the most promising configurations based on previous trial results.
    \item \textbf{Evolutionary Algorithms:} Useful for architectural search where the agent evolves the model structure through successive "generations" of candidates.
\end{itemize}

\subsection{Constraint-Driven Early Stopping}
One of the most powerful features of NNI integrated into this workflow is the \textbf{Assessor} module. In the context of Edge AI, a trial that exceeds the target MCU's Flash or RAM capacity is a "failed" candidate regardless of its accuracy. By integrating feedback from the STEdgeAI conversion tool, the NNI Assessor can implement early-stopping policies (e.g., Median Stopper or Curve Fitting) to terminate underperforming or non-compliant trials early. This significantly reduces the total optimization time, as the system can reallocate computational resources to candidates that are more likely to fit within the hardware's "Pareto front" of efficiency and performance.